\chapter{АЛГОРИТМ СИМУЛЯЦИИ БЕЗ ИСПОЛЬЗОВАНИЯ ГРАФИЧЕСКОГО ПРОЦЕССОРА}
\label{cha:cpu}

\section{Постановка задачи}

Ядром программного комплекса является решатель динамики твёрдых тел с
обнаружением столкновений и обработкой контактных ограничений. На данном этапе
реализован алгоритм симуляции, исполняемый на центральном процессоре (CPU). Он
служит эталонной реализацией: относительно неё на последующих этапах проверяется
корректность GPU-версии (раздел~\ref{cha:gpu}). Работа выполнена в репозитории
\texttt{dyphur} (C++20); CPU-исполнение обеспечивается бэкендом OpenMP единого
вычислительного слоя.

\section{Модель твёрдого тела}

Состояние твёрдого тела описывается вектором
\begin{equation}
  \mathbf{s} = \bigl(\mathbf{r},\, \mathbf{q},\, \mathbf{v},\,
  \boldsymbol{\omega}\bigr),
  \label{eq:state}
\end{equation}
где $\mathbf{r}$ --- положение центра масс, $\mathbf{q}$ --- единичный кватернион
ориентации, $\mathbf{v}$ --- линейная и $\boldsymbol{\omega}$ --- угловая
скорости~\cite{featherstone2008}. Уравнения движения Ньютона--Эйлера:
\begin{align}
  \dot{\mathbf{r}} &= \mathbf{v}, &
  m\dot{\mathbf{v}} &= \mathbf{F}, \\
  \dot{\mathbf{q}} &= \tfrac12\,\boldsymbol{\omega}\otimes\mathbf{q}, &
  \mathbf{I}\dot{\boldsymbol{\omega}} &=
  \boldsymbol{\tau} - \boldsymbol{\omega}\times(\mathbf{I}\boldsymbol{\omega}),
  \label{eq:newton-euler}
\end{align}
где $m$ --- масса, $\mathbf{I}$ --- тензор инерции, $\mathbf{F}$ и
$\boldsymbol{\tau}$ --- результирующие сила и момент, $\otimes$ ---
кватернионное произведение.

Данные тел хранятся в раскладке <<структура массивов>> (Structure of Arrays,
SoA): отдельные массивы для координат, скоростей, ориентаций и т.\,д. Такая
раскладка обеспечивает коалесцированный доступ к памяти и одинаково эффективна
как для векторизации на CPU, так и для исполнения на GPU.

\section{Метод численного интегрирования}

В качестве метода интегрирования выбран \emph{симплектический} (полунеявный)
метод Эйлера: сначала по силам обновляется скорость, затем по обновлённой
скорости --- положение:
\begin{align}
  \mathbf{v}_{n+1} &= \mathbf{v}_n + h\,\mathbf{M}^{-1}\mathbf{F}_n, \\
  \mathbf{r}_{n+1} &= \mathbf{r}_n + h\,\mathbf{v}_{n+1},
  \label{eq:symplectic}
\end{align}
где $h$ --- шаг интегрирования, $\mathbf{M}$ --- матрица масс. В отличие от явного
метода Эйлера, симплектический метод сохраняет фазовый объём и не накапливает
энергию, оставаясь устойчивым при типичных для реального времени шагах
($h \le 10$~мс). Ориентация интегрируется через производную кватерниона
\eqref{eq:newton-euler}~\cite{diebel2006} с последующей перенормировкой $\mathbf{q}$ для
компенсации дрейфа. Для повышения устойчивости шаг разбивается на $n$
подшагов (substepping), что эквивалентно уменьшению эффективного $h$.

\section{Обнаружение столкновений}

Обнаружение столкновений выполняется в две фазы.

\textbf{Широкая фаза} (broadphase) строит линейную иерархию ограничивающих
объёмов (LBVH)~\cite{karras2012}. Центры AABB тел кодируются 10-битными кодами Мортона, после чего
сортировка по ключу упорядочивает тела вдоль кривой, заполняющей пространство;
по отсортированному массиву строится двоичное дерево, обход которого даёт
кандидатные пары перекрывающихся AABB. Кодирование Мортона и построение LBVH
естественно параллелятся, что важно для GPU-версии.

\textbf{Узкая фаза} (narrowphase) для каждой кандидатной пары вычисляет точки
контакта и нормали методом разделяющей оси (Separating Axis Theorem, SAT).
Реализованы пары сфера--сфера, сфера--бокс и бокс--бокс. Результатом является
набор контактных точек с глубиной проникновения и нормалью.

\section{Решатель контактных ограничений}

Контактные и сочленённые ограничения разрешаются методом расширенной
позиционной динамики (Extended Position-Based Dynamics, XPBD)~\cite{macklin2016xpbd}.
Каждое ограничение $C(\mathbf{x})=0$ снабжается податливостью (compliance)
$\tilde{\alpha}=\alpha/h^2$; на каждом подшаге вычисляется приращение множителя
Лагранжа
\begin{equation}
  \Delta\lambda =
  \frac{-C - \tilde{\alpha}\,\lambda}
       {\nabla C\,\mathbf{M}^{-1}\nabla C^{\top} + \tilde{\alpha}},
  \qquad
  \Delta\mathbf{x} = \mathbf{M}^{-1}\nabla C^{\top}\,\Delta\lambda,
  \label{eq:xpbd}
\end{equation}
после чего позиции корректируются на $\Delta\mathbf{x}$. Ограничения
обходятся последовательно (схема Гаусса--Зейделя). XPBD устойчив, не зависит от
жёсткости в смысле явного шага и хорошо согласуется с подшаговым интегрированием.

\begin{lstlisting}[language=C++,caption={Шаг симуляции: интегрирование, столкновения, решатель},label=lst:step]
void step(World& w, float dt, int n_substeps) {
    const float h = dt / n_substeps;
    for (int s = 0; s < n_substeps; ++s) {
        integrate_velocities(w, h);   // v += h * M^-1 * F (симплектический Эйлер)
        predict_positions(w, h);      // x* = x + h * v
        auto pairs    = broadphase(w);          // LBVH -> кандидатные пары
        auto contacts = narrowphase(w, pairs);  // SAT -> контактные точки
        solve_xpbd(w, contacts, h);             // позиционные поправки
        update_velocities(w, h);      // v = (x - x_prev) / h
    }
}
\end{lstlisting}

\screenshot{0.85}{Симуляция падения и укладки твёрдых тел (CPU)}{Снимок экрана:
визуализация результата CPU-симуляции --- множество твёрдых тел, осевших на
опорную плоскость; кадр воспроизведения траектории.}{cpu-sim}

\section{Результаты}

Реализован полный цикл CPU-симуляции: симплектическое интегрирование с
подшагами, широкая фаза на основе LBVH, узкая фаза на основе SAT и решатель
ограничений XPBD. Реализация покрыта модульными тестами
(интегратор, широкая и узкая фазы, решатель, сочленения), а также сквозным
тестом устойчивости и детерминизма. CPU-версия используется как эталон
корректности для GPU-версии и как опорная точка для оценки производительности
(раздел~\ref{cha:test}).
